\section{Reduction to stuck graphs}\label{sec:reduction}

The number of connected graphs on $13$ vertices ($\num{50335907869219}\approx 5.0\times 10^{13}$) and on $14$ vertices ($\approx 2.9\times 10^{16}$; OEIS A001349) makes a naive search for the smallest member of $\I\setminus\IS$ impractical: deciding $\IS$ exactly costs microseconds to milliseconds per graph. The following observation turns the search for a smallest counterexample into a search for a much rarer kind of graph.

\begin{definition}\label{def:stuck}
A graph $G$ is \emph{stuck} if it is $\Ztwo$-acyclic and has no $\IS$-deletable vertex and no $\IS$-deletable edge, i.e.\ $\Nb(v)\notin\IS$ for every vertex $v$ and $\Nb(u,v)\notin\IS$ for every edge $uv$ (recall that $\emptyset\notin\IS$). A graph $A$ with a distinguished vertex $u$ is \emph{almost stuck at $u$} if it is $\Ztwo$-acyclic, has no $\IS$-deletable edge, and no vertex other than $u$ is $\IS$-deletable ($u$ itself may or may not be).
\end{definition}

A stuck graph is trivially outside $\IS$: it admits no strong contractible move at all. Conversely:

\begin{lemma}[Minimality implies stuck]\label{lem:stuck}
\begin{enumerate}[label=(\alph*)]
\item Let $n_0$ be the smallest order of a $\Ztwo$-acyclic graph not in $\IS$. Every $\Ztwo$-acyclic graph on $n_0$ vertices not in $\IS$ is obtained from a stuck graph on $n_0$ vertices by successively adding edges $uv$ with $\Nb(u,v)\in\IS$ (transformations (I4)); we say it lies in the \emph{$\mathrm{I}_4$-closure} of that stuck graph. Those with the minimum number of edges are themselves stuck.
\item Let $n_1$ be the smallest order of a graph in $\I\setminus\IS$. Every graph in $\I\setminus\IS$ on $n_1$ vertices lies in the $\mathrm{I}_4$-closure of a stuck graph in $\I\setminus\IS$ on $n_1$ vertices, which is $K_4$-free if the original graph is. Those with the minimum number of edges (in the whole class, or in the $K_4$-free class) are stuck.
\end{enumerate}
\end{lemma}

\begin{lemma}[Stuck graphs are $2$-connected]\label{lem:2conn}
Let $G$ be a stuck graph with a cut vertex $u$, so that $G=A\cup_u B$ with $A$, $B$ connected, $V(A)\cap V(B)=\{u\}$ and no edges between $A-u$ and $B-u$. Then $A$ and $B$ are almost stuck at $u$ and have at least three vertices. Consequently, if there is no almost stuck graph with at most $m$ vertices, every stuck graph with at most $2m$ vertices is $2$-connected.
\end{lemma}

\begin{lemma}[Degree bounds]\label{lem:degrees}
Let $G$ be a $2$-connected stuck graph on $n\ge 2$ vertices. Then every edge of $G$ lies in a triangle, every common neighbourhood $\Nb(u,v)$ is non-empty, $\delta(G)\ge 4$ and $\Delta(G)\le n-2$.
\end{lemma}

The proofs are elementary and are given in \ref{app:proofs}. To test stuckness during an enumeration one would like to replace the condition ``$\notin\IS$'' on neighbourhoods by the cheaper ``not $\Ztwo$-acyclic''. This is legitimate in the range of this paper:

\begin{lemma}\label{lem:Z2test}
Let $G$ be a $\Ztwo$-acyclic graph on $n\le 14$ vertices without a universal vertex. Then $G$ is stuck if and only if no open neighbourhood $\Nb(v)$ and no non-empty common neighbourhood $\Nb(u,v)$ is $\Ztwo$-acyclic. The same holds for ``almost stuck'' when $n\le 11$.
\end{lemma}
\begin{proof}
All neighbourhoods involved have at most $n-2\le 12$ vertices, and for graphs with at most $12$ vertices $\Ztwo$-acyclic is equivalent to $\IS$ (Table~\ref{tab:card}, Theorem~\ref{thm:A}).
\end{proof}

\begin{proposition}\label{prop:almost}
There is no almost stuck graph with at most $11$ vertices. Consequently every stuck graph with at most $22$ vertices is $2$-connected.
\end{proposition}
\begin{proof}
An almost stuck graph in which $u$ is not $\IS$-deletable is stuck, hence $\Ztwo$-acyclic and outside $\IS$; by Table~\ref{tab:card} there is no such graph with at most $12$ vertices. Algorithm~\ref{alg:stuck} in mode \emph{almost}, which by Lemma~\ref{lem:Z2test} detects exactly the $\Ztwo$-acyclic graphs with no $\IS$-deletable edge and exactly one $\IS$-deletable vertex, was run on all connected graphs with $3\le n\le 11$ vertices ($\num{1006700565}$ graphs for $n=11$) and found none. The second statement is Lemma~\ref{lem:2conn}.
\end{proof}

\begin{theorem}[Reduction]\label{thm:C}
\begin{enumerate}[label=(\roman*)]
\item If $n_0\le 22$ is the smallest order of a $\Ztwo$-acyclic graph outside $\IS$, every such graph on $n_0$ vertices lies in the $\mathrm{I}_4$-closure of a stuck graph on $n_0$ vertices.
\item If $n_1\le 22$ is the smallest order of a graph in $\I\setminus\IS$, every graph in $\I\setminus\IS$ on $n_1$ vertices lies in the $\mathrm{I}_4$-closure of a stuck graph in $\I\setminus\IS$ on $n_1$ vertices, $K_4$-free if the graph is.
\end{enumerate}
In both cases every stuck graph on $n_i$ vertices is $2$-connected with $4\le\delta\le\Delta\le n_i-2$, and for $n_i\le 14$ stuckness is decided by $\Ztwo$-homology of the neighbourhoods (Lemma~\ref{lem:Z2test}). Hence, to find all graphs of either kind on $n_i$ vertices it suffices to enumerate the $2$-connected graphs with these degree bounds, keep the stuck ones, and compute their $\mathrm{I}_4$-closures.
\end{theorem}

For $n=13$ the enumeration of $2$-connected graphs with $4\le\delta$ and $\Delta\le 11$ comprises $\num{21011415992470}$ graphs, about $2.4$ times fewer than all connected graphs; the real gain lies in the stuck filter: the cascade of Section~\ref{sec:algorithms} rejects all but a fraction $7.5\times10^{-4}$ of the graphs by looking at cones, the Euler characteristic and a few neighbourhoods, computes the homology of $G$ itself for only $\num{15693920490}$ graphs, and never runs the exact $\IS$ decision during the enumeration. For $n=14$ we restricted the search to $K_4$-free graphs (clique complexes of dimension $2$); this class contains the dunce hat $\Dunce$ and comprises $\num{64856470381820}$ graphs under the same degree constraints. By Theorem~\ref{thm:C}(ii) this determines all $K_4$-free graphs in $\I\setminus\IS$ on $14$ vertices.
